\hspace{3mm}

\centerline{\bf \Huge Разнобой 8 класс}

\begin{flushright}
        \scriptsize{ Одного было бы мало! \\ Левая голова Огр мага}
\end{flushright}

\problem{1}
На доске написаны четыре числа, ни одно из которых не равно $0$. Если каждое из них умножить на сумму трёх остальных, получатся четыре одинаковых результата. Докажите, что квадраты записанных на доске чисел равны.

\problem{2}
Три спортсмена пробежали дистанцию в $3$ километра. Первый километр они бежали с постоянными скоростями $v_1, v_2$ и $v_3$ соответственно, такими, что $v_1 > v_2 > v_3$. После отметки в $1$ километр каждый из них изменил скорость: первый $-$ с $v_1$ на $v_2$, второй $-$ с $v_2$ на $v_3$, а третий $-$ с $v_3$ на $v_1$. Кто из спортсменов пришел к финишу последним?

\problem{3}
Даны два ненулевых числа. Если к каждому из них прибавить единицу, а также из каждого из них вычесть единицу, то сумма обратных величин четырёх полученных чисел будет равна 0. Какое число может получиться, если из суммы исходных чисел вычесть сумму их обратных величин? Найдите все возможности.

\problem{4}
Будем называть две клетки клетчатой таблицы соседями, если у них есть общая сторона. Можно ли покрасить в белой таблице размером $10 \times 10$ клеток $32$ клетки в черный цвет так, чтобы у каждой черной клетки было поровну черных и белых соседей, а у каждой белой клетки $-$ не поровну?


\newpage

\hspace{3mm}

\centerline{\bf \Huge Разнобой 9 класс}

\begin{flushright}
        \scriptsize{Но мы всё равно не умеем считать!\\ Правая голова Огр мага
}
\end{flushright}
 
\problem{1} 
Два приведённых квадратных трёхчлена $f(x)$ и $g(x)$ таковы, что каждый из них имеет по два корня, и выполняются равенства $f(1) = g(2)$ и $g(1) = f(2)$. Найдите сумму всех четырёх корней этих трёхченов.

\problem{2}
Петя и Миша стартуют по круговой дорожке из одной точки в направлении против часовой стрелки. Оба бегут с постоянными скоростями, скорость Миши на 2\% больше скорости Пети. Петя всё время бежит против часовой стрелки, а Миша может менять направление бега в любой момент, непосредственно перед которым он пробежал полуокружность или больше в одном направлении. Покажите, что пока Петя бежит первый круг, Миша может трижды, не считая момента старта, поравняться (встретиться или догнать) с ним.

\problem{3}
Вася записал в клетки таблицы $9 \times 9$ натуральные числа от $1$ до $81$ (в каждой клетке стоит по числу, все числа различны). Оказалось, что любые два числа, отличающихся на $3$, стоят в соседних по стороне клетках. Верно ли, что обязательно найдутся две угловых клетки, разность чисел в которых делится на $6$?

\problem{4}
На доске девять раз (друг под другом) написали некоторое натуральное число $N$. Петя к каждому из $9$ чисел приписал слева или справа одну ненулевую цифру; при этом все приписанные цифры различны. Какое наибольшее количество простых чисел могло оказаться среди $9$ полученных чисел?